\section{Problem Statement}
\label{sec:problem}

\paragraph{Prediction and residuals.}
An image $x$ has a finite set $\cP$ of pixels with valid ground-truth depth $y(p)>0$.
A frozen monocular metric model $f$ predicts $\yhat(p)=f(x)(p)$; nothing is retrained.
We work with the log-depth residual
\begin{equation}
d(p)=\log y(p)-\log\yhat(p),
\label{eq:residual}
\end{equation}
so an interval that is symmetric in $d$ is multiplicative in depth: $C(p)=[\yhat(p)e^{m(p)-w(p)},\,\yhat(p)e^{m(p)+w(p)}]$ for a centre shift $m$ and a half-width $w$.

\paragraph{A primer on conformal prediction.}
Let $R=\mathcal{R}(x,y)$ be a nonconformity score, \eg $|d|/s$ for a scale map $s$, computed on $N$ calibration units and on one test unit.
If the $N+1$ scores are exchangeable (their joint law is invariant under permutations; i.i.d.\ is a special case), the test score is equally likely to take any rank among them.
Hence the $\lceil(N+1)(1-\alpha)\rceil$-th smallest calibration score $\qhat$ satisfies $\P(R_{\text{test}}\le\qhat)\ge1-\alpha$.
The interval $\{y:\mathcal{R}(x,y)\le\qhat\}$ therefore covers the truth with probability at least $1-\alpha$, for any model and any data distribution.
This is \emph{split conformal prediction}; its only assumption is exchangeability.
For dense outputs, conformal risk control~\cite{angelopoulos2024crc} chooses $\qhat$ so that the expected value of a bounded loss, such as the fraction of a new image's pixels left uncovered, is at most $\alpha$.

\paragraph{Why pixels are not the unit.}
If a model misjudges the size of a room by 20\%, every pixel in that image is wrong together.
Pixels of one image are therefore far from independent, and treating them as exchangeable calibration points gives no valid statement.
The exchangeable unit is the \emph{image}, or the \emph{scene} (room, drive, visit) when several images share one.
A guarantee built this way holds \emph{on average over images}.
Individual images can still be covered anywhere from 40\% to 100\%, because each image has its own scale error.

\paragraph{Sparse metric measurements.}
A sensor returns $k$ measurements per image.
At a measured pixel it reports a depth $\ytil(p)$ and a residual $\dtil(p)=\log\ytil(p)-\log\yhat(p)$.
For a perfect sensor $\dtil=d$; for a real one $\dtil(p)=d(p)+e(p)$ with a sensor error $e$.
The positions follow a \emph{layout}, a random rule that may depend on $x$ and $\yhat$ but never on $y$.
We use five: \emph{uniform} (a uniform sample of $\cP$), \emph{rows} (pixels on four random image rows, a few-beam scanner), \emph{cluster} (the $k$ pixels nearest one random centre, a single observed patch), \emph{edges} (uniform among edge pixels of $\yhat$, like feature points) and \emph{near} (uniform among the 30\% of pixels with the smallest $\yhat$, a range-limited sensor).
In the main method, $k-n$ measurements follow the layout and form the \emph{structured anchor} set $\cS$.
The remaining $n$ form the \emph{probe} set $\cQ$, drawn uniformly without replacement from $\cU=\cP\setminus\cS$.
Intervals are evaluated on $\cE=\cU\setminus\cQ$.
For a pixel $p$, $\delta(p)$ is its distance to the nearest measurement, divided by the image diagonal.
Random probes need a sensor that can be read at random pixels: a dense or semi-dense sensor subsampled in software (phone LiDAR, time-of-flight or structured-light cameras), or a steerable scanner~\cite{pittaluga2020mems,bergman2020adaptivelidar}.

\paragraph{Coverage targets.}
A method outputs $C(p)=[\ell(p),u(p)]$ at every evaluation pixel.
The coverage of image $i$ is
\begin{equation}
\Cov_i=\frac{1}{|\cE_i|}\sum_{p\in\cE_i}\1\{y(p)\in C(p)\}.
\label{eq:cov}
\end{equation}
At level $1-\alpha$ we distinguish the following targets:
\begin{itemize}[leftmargin=*,itemsep=1pt,topsep=2pt]
\item \emph{Marginal (image or scene):} $\E[\Cov_i]\ge1-\alpha$ over images (or scenes), measurements and calibration data.
\item \emph{Image event:} with probability at least $1-\alpha$, a new image has at least a fraction $1-\beta$ of its pixels covered.
\item \emph{Per-image:} $\E[\Cov_i\mid x_i,\cS_i,\dtil|_{\cS_i}]\ge1-\alpha$ for \emph{every} image, with only the probes random.
No image is then always under-covered.
\item \emph{Per distance bin:} for bins of $\delta$ (edges $0.02,0.05,0.1,0.2$), the expected coverage of the pixels of a new scene in each bin is at least $1-\alpha$, so far pixels are not traded for near ones.
\end{itemize}
Every guarantee concerns the values used to calibrate: sensor readings certify the readings, not the depth.
We always \emph{evaluate} against independent ground truth $y$.

\paragraph{Efficiency.}
Among valid methods we prefer narrower intervals.
We report the median relative width $(u-\ell)/\yhat$ and the interval score $(u-\ell)+\frac{2}{\alpha}[(\ell-y)_++(y-u)_+]$, which also penalizes misses by their size.
The \emph{relative interval score} (IS) divides the interval score by $y$, averages over the pixels of an image, and takes the median over images.
It keeps the miss penalty while staying robust to a few extreme images.

\paragraph{Research questions.}
(Q1)~Which unit makes conformal guarantees for dense depth valid, and what do they cost?
(Q2)~How much do sparse metric anchors help, and do their guarantees survive realistic layouts, layout mismatch and sensor noise?
(Q3)~Can a per-image guarantee be kept for \emph{any} layout and any centre, and what does the sensor budget buy?
(Q4)~Which centre should be calibrated, and what exactly is certified when the measurements come from a real, imperfect sensor?
